\documentclass[12pt,a4paper,oneside]{article}
\usepackage[brazil]{babel}
\usepackage[utf8]{inputenc}
\usepackage{olymp}
\usepackage{color}
\usepackage[dvipsnames]{xcolor}
\usepackage{listings}
\usepackage{subfigure}
\setcounter{problem}{4}

\begin{document}

\begin{problem}{2048}{2048}{2}
 
Há poucos anos um jogo se popularizou pela internet, o 2048. A ideia do jogo é bem simples. Utiliza-se sempre um tabuleiro $4\times 4$, onde cada célula poderá estar vazia ou conter uma peça com um valor inteiro, sempre uma potência de 2 entre 2 e 1024.

O jogador pode então escolher entre quatro tipos de movimentos (para cima, para baixo, para a esquerda ou para a direita), desde que seja possível realizar tais movimentos. Ao realizar um movimento, todas as peças são movimentadas o máximo possível na direção escolhida. Além disso, se um movimento forçar duas peças de igual valor uma contra a outra, as duas peças se ``fundem'', tendo seus valores somados para obter o valor da peça resultante. Após a realização de um movimento, uma nova peça de valor 2 ou 4 surge aleatoriamente em uma das células vazias do tabuleiro. 

As Figuras 1 e 2 ilustram um movimento para cima. A Figura 1 mostra a configuração inicial do tabuleiro, enquanto a Figura 2 ilustra a configuração do mesmo tabuleiro após o jogador apertar a tecla para cima. Repare que, além de levar todas as peças o máximo possível para a parte de cima do tabuleiro, o movimento fez com que as duas peças de valor 2 na segunda coluna se fundissem e formassem uma nova peça de valor 4. Repare também que, após o movimento, uma nova peça de valor 2 surgiu em uma das posições vazias do tabuleiro (neste caso, na posição logo abaixo da peça de valor 8).

    \begin{figure}[!h]
    	\centering
    	\begin{minipage}[t]{.45\textwidth}
    		\centering
    		\includegraphics[scale=0.25]{images/exercise_70_0.png}
    		\caption{Configuração inicial}
    	\end{minipage}
    	\begin{minipage}[t]{.45\textwidth}
    		\centering
    		\includegraphics[scale=0.25]{images/exercise_70_1.png}
    		\caption{Configuração após um movimento para cima.}
    	\end{minipage}
    \end{figure}

% \begin{figure}[h] 
%   \subfigure[Picture 1]{% 
%     \includegraphics[scale=0.25]{images/exercise_70_0.png} \label{fig:1} 
%   } 
%   \quad 
%   \subfigure[Picture 2]{% 
%     \includegraphics[scale=0.25]{images/exercise_70_1.png} \label{fig:2} 
%   } 
%   \caption{Main figure caption} 
% \end{figure}

% \begin{figure}[h]
%   \begin{subfigure}[b]{0.5\textwidth}
%     \includegraphics[scale=0.25]{images/exercise_70_0.png}
%     \caption{Picture 1}
%     \label{fig:1}
%   \end{subfigure}
%   %
%   \begin{subfigure}[b]{0.5\textwidth}
%     \includegraphics[scale=0.25]{images/exercise_70_0.png}
%     \caption{Picture 2}
%     \label{fig:2}
%   \end{subfigure}
% \end{figure}

% \begin{minipage}{0.5\textwidth}
% \begin{figure}
%     \centering
%     \includegraphics[scale=0.25]{images/exercise_70_0.png}
%     \caption{Configuração inicial}
%     \label{fig:my_label}
% \end{figure}
% \end{minipage}
% \begin{minipage}{0.5\textwidth}
% \begin{figure}
%     \centering
%     \includegraphics[scale=0.25]{images/exercise_70_1.png}
%     \caption{Configuração após um movimento para cima.}
%     \label{fig:my_label}
% \end{figure}
% \end{minipage}

O objetivo do jogo é realizar diversos movimentos consecutivos até obter uma peça de valor 2048. O problema é que existem diversas configurações de tabuleiro para as quais não existe nenhum movimento possível. Ao alcançar uma dessas configurações, o jogo acaba (\textit{Game over!} - veja a Figura 3). É importante observar que mesmo que não exista nenhuma célula vazia, isto não configura necessariamente uma situação de \textit{Game over}, uma vez que ainda é possível haver movimentos quando há peças de igual valor adjacentes uma a outra (veja a Figura 4).

    \begin{figure}[!h]
    	\centering
    	\begin{minipage}[t]{.45\textwidth}
    		\centering
    		\includegraphics[scale=0.25]{images/exercise_70_2.png}
    		\caption{Situação de \textit{Game over!}}
    	\end{minipage}
    	\begin{minipage}[t]{.45\textwidth}
    		\centering
    		\includegraphics[scale=0.25]{images/exercise_70_3.png}
    		\caption{Ainda há movimentos disponíveis}
    	\end{minipage}
    \end{figure}

Faça um programa para analisar uma configuração de tabuleiro do jogo 2048 e determinar se existe ou não algum movimento possível. Ou seja, seu programa deve determinar se a configuração do tabuleiro indica uma situação de \textit{Game over!}.

%\Notes
%
%Observações, se tiver.

\newpage
\InputFile

O programa terá apenas um caso de teste.

O caso de teste é composto por quatro linhas com quatro valores inteiros cada, indicando a configuração do tabuleiro a ser analisado. Valores $0$ indicam uma célula vazia, enquanto os outros valores serão sempre potências de $2$ indicando o valor da peça contida em cada célula (entre $2$ e $1024$, inclusive).

\OutputFile

Seu programa deve gerar apenas uma linha de saída, contendo uma das frases a seguir:
\begin{itemize}
 \item ``\texttt{Keep playing!}'', caso ainda haja movimentos possíveis;
 \item ``\texttt{Game over!}'', caso não haja mais movimentos possíveis.
 \end{itemize}
 

%\Notes
%Nesta questão, seu programa deve usar a função \texttt{fatorial} dada %acima exatamente como está, não a modifique! Sua
%tarefa é apenas criar o programa com a função \texttt{main}.

\Example

\begin{example}
\exmp{
0 2 0 0
0 0 0 0
0 0 2 4
0 2 8 2
}{
Keep playing!
}%
\end{example}

\begin{example}
\exmp{
2 4 32 4
4 16 128 16
8 32 2 4
16 2 4 32
}{
Game over!
}%
\end{example}

\begin{example}
\exmp{
2 4 32 4
4 2 128 16
2 4 16 2
2 2 4 4
}{
Keep playing!
}%
\end{example}

%Comentários sobre o exemplo, se necessário.

\end{problem}

\end{document}
